%% The order-topological duality for Heyting algebras.

\section{Esakia spaces and Heyting homomorphisms}

\begin{annexitem}{annex:p10.esakia.1}{thm:p10-esakia-objects}{Heyting algebras and Esakia spaces}
\begin{proof}
By Theorem~\ref{thm:p10-heyting-prime-representation}, a prime filter
\(P\) omits \(a\Rightarrow b\) exactly when a prime extension
\(Q\supseteq P\) contains \(a\) and omits \(b\).
The latter condition says precisely
\[
 P\in\downarrow(\widehat a\setminus\widehat b).
\]
This proves the displayed identity in the theorem.

The Priestley topology on \(X_H\) has a basis of sets
\(\widehat a\setminus\widehat b\).
Every clopen \(C\) is a finite union of such basic sets:
cover \(C\) by all basic sets contained in it, add its open
complement, and apply compactness.
The downward closure of each basic set is clopen by the
displayed identity.
Downward closure commutes with unions, so \(\downarrow C\)
is a finite union of clopens and is clopen.
Thus \(X_H\) is Esakia.

Conversely, suppose \(X\) is Esakia and \(U,V\) are clopen upsets.
The difference \(U\setminus V\) is clopen, so
\(W=X\setminus\downarrow(U\setminus V)\) is clopen.
It is an upset because the complement of a downset is an upset.
For any upset \(T\),
\[
 T\subseteq W\quad\Longleftrightarrow\quad T\cap U\subseteq V.
\]
For the forward direction, if \(x\in T\cap U\) and \(x\notin V\),
then \(x\in\downarrow(U\setminus V)\), contradicting \(x\in W\).
For the reverse direction, if \(x\in T\setminus W\), choose
\(y\geq x\) with \(y\in U\setminus V\).
Since \(T\) is an upset, \(y\in T\), contradicting
\(T\cap U\subseteq V\).
This residuation law, in particular for clopen upsets \(T\),
proves the asserted Heyting operation on \(\operatorname{ClopUp}(X)\).

Priestley representation makes \(\eta_H\) a bounded lattice
isomorphism onto all clopen upsets.
The implication identity proved above makes it preserve
\(\Rightarrow\), so it is a Heyting isomorphism.
Priestley reconstruction supplies the evaluation order
isomorphism and homeomorphism for \(X\).
For a trivial Heyting algebra the prime-filter space is empty;
all displayed subsets are empty, and its clopen-upset algebra
is trivial. Thus these conclusions include both edge cases.
\end{proof}
\end{annexitem}

\begin{annexitem}{annex:p10.esakia.2}{thm:p10-heyting-prime-lift}{Lifting a prime filter}
\begin{proof}
The complement \(H\setminus G\) is a lattice ideal, since \(G\)
is a prime filter.
In \(K\), define
\[
\begin{aligned}
 E&=\{k:\exists f\in F\ \exists g\in G\
                              f\wedge h(g)\leq k\},\\
 J&=\{k:\exists b\in H\setminus G\quad k\leq h(b)\}.
\end{aligned}
\]
The set \(E\) is the filter generated by \(F\cup h[G]\).
It contains \(1\), is upward closed, and is meet-closed by
combining witnesses \(f,g\) by their respective meets.
It contains \(F\) by taking \(g=1\), and \(h[G]\) by taking \(f=1\).
The set \(J\) is the ideal generated by \(h[H\setminus G]\):
it contains \(0\), is downward closed, and is join-closed
because \(H\setminus G\) is an ideal and \(h\) preserves joins.

They are disjoint.
Otherwise some \(f\in F\), \(g\in G\), and \(b\notin G\)
would satisfy
\[
 f\wedge h(g)\leq h(b).
\]
Residuation in \(K\) and preservation of implication by \(h\) give
\[
 f\leq h(g)\Rightarrow h(b)=h(g\Rightarrow b).
\]
Upward closure of \(F\) yields
\(g\Rightarrow b\in h^{-1}[F]\subseteq G\).
Since \(g\in G\) and \(g\wedge(g\Rightarrow b)\leq b\),
the filter \(G\) must contain \(b\), a contradiction.

Prime-filter separation in \(K\) now gives a prime filter
\(F'\supseteq E\) avoiding \(J\).
It contains \(F\).
For every \(g\in G\), \(h(g)\in E\subseteq F'\);
for every \(b\notin G\), \(h(b)\in J\) and is absent from \(F'\).
Hence \(h^{-1}[F']=G\), as required.
Only finitely many witnesses are used in each disjointness
argument; the prime extension is the explicitly assumed BPIT step.
\end{proof}
\end{annexitem}

\begin{annexitem}{annex:p10.esakia.3}{thm:p10-esakia-functors}{Functorial Esakia duality}
\begin{proof}
For a Heyting homomorphism \(h:H\to K\), Priestley duality
already makes the inverse-image map continuous and
order-preserving on prime filters.
The preceding lifting lemma supplies its back condition.
Thus it is an Esakia morphism.

Conversely, let \(g:X\to Y\) be an Esakia morphism.
Inverse image preserves the bounds, finite meets and joins,
and sends clopen upsets to clopen upsets.
It also preserves implication.
To see this, for clopen upsets \(U,V\subseteq Y\) and \(x\in X\),
suppose \(g(x)\in U\Rightarrow V\).
For every \(z\geq x\), monotonicity gives \(g(z)\geq g(x)\);
if \(g(z)\in U\), it follows that \(g(z)\in V\).
Thus \(x\in g^{-1}[U]\Rightarrow g^{-1}[V]\).

If instead \(g(x)\notin U\Rightarrow V\), choose
\(y\geq g(x)\) in \(U\setminus V\).
Back supplies \(z\geq x\) with \(g(z)=y\).
This witnesses
\(x\notin g^{-1}[U]\Rightarrow g^{-1}[V]\).
We have proved
\[
 g^{-1}[U\Rightarrow V]
       =g^{-1}[U]\Rightarrow g^{-1}[V].
\]
Hence the inverse-image map is a Heyting homomorphism.

Identity maps satisfy back, and composites satisfy it by
lifting first through the second map and then through the first.
The continuous monotone conditions and inverse-image laws
also compose, so the constructions are contravariant functors.
The object representation maps are Heyting isomorphisms by
the preceding object theorem.
The evaluation order isomorphisms are Esakia isomorphisms:
an order isomorphism has back, as does its inverse, and both
directions are continuous.

The same membership equations as in Priestley duality prove
naturality and the identity reconstruction composites.
For \(h:H\to K\), the preimage of \(\widehat a\) under its
prime-filter map is \(\widehat{h(a)}\).
For \(g:X\to Y\), pulling the point filter at \(x\) back
along the induced algebra map gives the point filter at \(g(x)\).
Since the maps now preserve all the additional Heyting and
Esakia structure, these identities establish the stated
contravariant equivalence.
All pointwise conditions are vacuous on an empty domain,
and the trivial-algebra cases agree with the earlier Priestley
correspondence.
\end{proof}
\end{annexitem}
